% =====================================================================
% Report content to be pasted INSIDE \begin{document} ... \end{document}
%
% Required packages in your Overleaf preamble:
%   \usepackage{graphicx}
%   \usepackage{amsmath, amssymb}
%   \usepackage{float}        % for [H] figure placement
%   \usepackage{booktabs}     % nice tables (\toprule, \midrule, \bottomrule)
%   \usepackage{listings}     % code snippets
%   \usepackage{xcolor}
%   \usepackage{hyperref}     % cross-references
%   \lstset{basicstyle=\ttfamily\footnotesize, breaklines=true,
%           keywordstyle=\color{blue}, commentstyle=\color{gray!80!black},
%           numbers=left, numberstyle=\tiny, stepnumber=1, frame=single,
%           language=Matlab}
%
% Place the images folder next to your main.tex
% =====================================================================

\title{Quarter-Car Modeling and Control \\ \large Project 1 -- Driver Assistance System Design}
\author{Dennis [Surname] \\ Politecnico di Torino}
\date{\today}
\maketitle

%======================================================================
\section{Introduction}
%======================================================================
The vertical dynamics of a road vehicle dominates ride comfort, road
holding and tyre wear. The classical \emph{quarter-car} model -- two
masses linked by a spring--damper element and connected to the ground
through the tyre stiffness and damping -- captures the essential
phenomena with only two degrees of freedom, while remaining simple
enough to support analytical insight and controller design.

This project investigates the quarter-car model in three progressive
steps:
\begin{itemize}
    \item \textbf{Exercise 1} -- frequency-domain analysis of the
          linear model under passive, skyhook and groundhook
          configurations;
    \item \textbf{Exercise 2} -- time-domain simulation under random
          road excitation, and implementation of a
          \emph{real-world} skyhook law on a semi-active suspension;
    \item \textbf{Exercise 3} -- nonlinear semi-active suspension with
          a current-controlled magneto-rheological damper, parameter
          tuning under bump and random profiles, and Bode analysis of
          the linearized system around two extreme operating points.
\end{itemize}

The same nominal vehicle parameters are used throughout, summarised in
Table~\ref{tab:params}. The damping coefficient $c$ in the actual code
is $1500\,\mathrm{Ns/m}$, slightly larger than the value suggested in
the project assignment in order to model a stiffer reference passive
damper representative of a medium-sized passenger car.

\begin{table}[H]
    \centering
    \begin{tabular}{lcl}
        \toprule
        \textbf{Symbol} & \textbf{Value} & \textbf{Description} \\
        \midrule
        $m_s$  & $400$\,kg     & Sprung mass (chassis) \\
        $m_u$  & $50$\,kg      & Unsprung mass (wheel + suspension) \\
        $K$    & $25\,000$\,N/m   & Suspension stiffness \\
        $P$    & $150\,000$\,N/m  & Tyre vertical stiffness \\
        $c$    & $1500$\,Ns/m   & Passive damping coefficient \\
        $c_t$  & $0$\,Ns/m      & Tyre damping coefficient \\
        \bottomrule
    \end{tabular}
    \caption{Quarter-car nominal parameters.}
    \label{tab:params}
\end{table}

%======================================================================
\section{Quarter-Car Model -- Theoretical Background}
\label{sec:theory}
%======================================================================
The 2-DOF passive quarter-car equations of motion are
\begin{align}
    m_s \ddot z_s &+ c(\dot z_s - \dot z_u) + K(z_s - z_u) = 0,
    \label{eq:passive_s}\\
    m_u \ddot z_u &+ c(\dot z_u - \dot z_s) + c_t(\dot z_u - \dot r) +
    K(z_u - z_s) + P(z_u - r) = 0,
    \label{eq:passive_u}
\end{align}
where $z_s$ and $z_u$ are the vertical displacements of the sprung and
unsprung masses and $r$ is the road profile.

\paragraph{Skyhook and groundhook.}
An \emph{ideal} skyhook damper $c_s$ would connect the sprung mass to
an inertial reference above the vehicle, removing low-frequency
oscillations of the chassis. An \emph{ideal} groundhook damper $c_g$
would connect the unsprung mass to the ground, reducing tyre-force
fluctuations and improving handling. Both are physically unfeasible,
but the same dynamic effect can be reproduced by an actuator placed
between the two masses producing the force
\begin{equation}
    F = -(c_s + c)\dot z_s + (c + c_g)\dot z_u .
    \label{eq:realworld}
\end{equation}

\paragraph{Damping ratios.}
Defining the critical damping
$c_{cr} = 2\sqrt{K m_s}$, all damping coefficients are normalised:
\begin{equation}
    \zeta   = \tfrac{c}{c_{cr}}, \quad
    \zeta_t = \tfrac{c_t}{c_{cr}}, \quad
    \zeta_s = \tfrac{c_s}{c_{cr}}, \quad
    \zeta_g = \tfrac{c_g}{c_{cr}} .
\end{equation}
The control logic is conveniently expressed through two parameters,
\begin{equation}
    \zeta_m = \tfrac{\zeta_s + \zeta_g}{2}, \qquad
    \zeta_d = \tfrac{\zeta_s - \zeta_g}{2}, \qquad
    \zeta_0 = \zeta + \zeta_m ,
\end{equation}
where $\zeta_d > 0$ corresponds to skyhook behaviour and
$\zeta_d < 0$ to groundhook behaviour.

\paragraph{State-space representation.}
Selecting $\mathbf{x} = [z_s, \dot z_s, z_u, \dot z_u, r]^{T}$ and
input $u = \dot r$, the system becomes
\begin{equation}
    \dot{\mathbf{x}} = A\mathbf{x} + B u, \quad
    \mathbf{y} = C\mathbf{x} + D u,
\end{equation}
with $\mathbf{y} = [\ddot z_s,\, F_t,\, F_{act}]^{T}$
(comfort, handling, actuator force). The road state $r$ is appended
in order to use $\dot r$ as a true input; this introduces a nil
eigenvalue in $A$. The matrices used in the MATLAB scripts are
generated symbolically in \texttt{symbolic\_state\_space\_gen.m}.

%======================================================================
\section{Exercise 1 -- Frequency-Response Analysis}
%======================================================================

\subsection{Theory}
Three transfer functions are required by the assignment:
\begin{equation}
    \frac{\ddot Z_s}{R}, \qquad
    \frac{F_t}{P R}, \qquad
    \frac{F_{act}}{P R} .
\end{equation}
Since the input of the state-space model is $\dot r$, MATLAB's
\texttt{bode} computes $X/\dot R$. Because $\dot R(j\omega) =
j\omega\,R(j\omega)$, the conversion
\begin{equation}
    \left|\frac{X}{R}\right| = \omega \cdot \left|\frac{X}{\dot R}\right|
\end{equation}
is required, and the second and third transfer functions are further
normalised by $P$. Each working condition is also analysed via the
eigenvalues of $A$, which describe damping ratios and natural
frequencies of the closed-loop system.

\subsection{Implementation}
Three damping configurations are looped over:
\begin{lstlisting}
switch damp
    case 1   % passive damper
        zeta0  = zeta;     zeta_d = 0;
    case 2   % skyhook
        zeta0  = 0.43;     zeta_d = 0.25;
    case 3   % groundhook
        zeta0  = 0.43;     zeta_d = -0.25;
end
\end{lstlisting}
For each case the $A,B,C,D$ matrices are built according to
Section~\ref{sec:theory} and a frequency vector
$f = \mathrm{linspace}(10^{-3},100,1000)$ rad/s is fed to
\texttt{bode}. The magnitudes returned are multiplied by
$\omega$ (and divided by $P$ for the second and third outputs):
\begin{lstlisting}
sys = ss(A,B,C,D);
[mag, phase, wout] = bode(sys, f);
if n == 1
    plot(wout, squeeze(mag(n,:,:)) .* wout)        % zddot_s / r
else
    plot(wout, squeeze(mag(n,:,:)) .* wout / P)    % F / (P*r)
end
\end{lstlisting}
The eigenvalues of $A$ are then computed by \texttt{eig(A)} and plotted
on the complex plane.

\subsection{Results}

\begin{figure}[H]
    \centering
    \includegraphics[width=\textwidth]{images/ex1_comparison.png}
    \caption{Frequency response of the quarter-car model under passive
    damper, skyhook and groundhook configurations: sprung mass
    acceleration $\ddot z_s/r$ (top), normalised tyre-force variation
    $F_t/(Pr)$ (middle) and normalised actuator force $F_{act}/(Pr)$
    (bottom).}
    \label{fig:ex1_comparison}
\end{figure}

Figure~\ref{fig:ex1_comparison} highlights the three classical
trade-offs:
\begin{itemize}
    \item Around the \emph{sprung mass} natural frequency
          ($\approx 7$\,rad/s) the skyhook strongly attenuates the
          peak of $\ddot z_s/r$, while the groundhook amplifies it.
          This confirms the comfort-oriented nature of skyhook
          control.
    \item Around the \emph{unsprung mass} resonance
          ($\approx 55$\,rad/s) the groundhook flattens the
          tyre-force peak, while the skyhook keeps it almost as large
          as in the passive case -- a clear handling improvement of
          groundhook.
    \item The actuator force is null in the passive case and exhibits
          a peak at one of the two resonances depending on which
          working condition is chosen (sprung resonance for skyhook,
          unsprung resonance for groundhook), giving a quick estimate
          of the actuator authority required.
\end{itemize}

\begin{figure}[H]
    \centering
    \includegraphics[width=0.85\textwidth]{images/ex1_eig.png}
    \caption{Eigenvalues (pole map) of the linearised quarter-car
    under the three working conditions.}
    \label{fig:ex1_eig}
\end{figure}

The pole map in Figure~\ref{fig:ex1_eig} shows two pairs of complex
poles for the passive damper (one for the chassis mode at $\sim
7$\,rad/s, one for the wheel mode at $\sim 56$\,rad/s). The
\textbf{skyhook} configuration shifts the chassis poles further into
the left half plane (more damping on $\ddot z_s$) while leaving the
wheel poles almost unchanged. The \textbf{groundhook}, conversely,
pushes the wheel poles deep to the left ($\sigma \approx -44$) and
relaxes the chassis poles, which become almost critically damped on
the real axis. The nil pole at the origin in each case comes from the
augmented road state $r$.

%======================================================================
\section{Exercise 2.1 -- Road-Profile Excitation}
%======================================================================

\subsection{Theory}
The road profile is now a stochastic disturbance representative of a
real surface. The model is fed with band-limited white noise filtered
to obtain a frequency content consistent with vehicle speed: the
spatial PSD of the road becomes a temporal PSD that scales with
$V$ (the longitudinal speed). The actuator force-velocity diagram
\begin{equation}
    F_{act} \quad \text{vs} \quad v_{rel} = \dot z_s - \dot z_u
\end{equation}
characterises the type of damper: a linear viscous damper would lie
on a straight line, while skyhook/groundhook produce dispersed clouds
because the force depends on absolute velocities rather than on
the relative one.

\subsection{Implementation -- Simulink model}

\begin{figure}[H]
    \centering
    \includegraphics[width=0.9\textwidth]{images/sim2_1.png}
    \caption{Simulink model \texttt{Exercise2SIM.slx} used in
    Exercise 2.1.}
    \label{fig:sim2_1}
\end{figure}

The Simulink model in Figure~\ref{fig:sim2_1} feeds a
\emph{Band-Limited White Noise} block through the road shaping filter
\begin{equation}
    G_r(s) = \frac{2\pi\sqrt{6.4\!\cdot\!10^{-7}\,V}}{s + 1.22},
\end{equation}
whose gain scales with the vehicle speed $V$. The derivative block
produces $\dot r$, which is the input of the state-space block
implementing the $(A,B,C,D)$ matrices written by the MATLAB script.
The outputs are demultiplexed: the relative velocity is computed as
$\dot z_s - \dot z_u$ and logged together with the actuator force
through \texttt{V\_act} and \texttt{F\_act} \texttt{To-Workspace}
blocks. The MATLAB driver sweeps the three vehicle speeds and the
three damping configurations:
\begin{lstlisting}
speeds = [10 30 50];
for n = 1:length(speeds)
    V = speeds(n);
    for damp = 1:3
        % set zeta0, zeta_d, build A,B,C,D
        out = sim("Exercise2SIM.slx");
        plot(out.V_act, out.F_act, '.')
    end
end
\end{lstlisting}

\subsection{Results}

\begin{figure}[H]
    \centering
    \includegraphics[width=0.85\textwidth]{images/ex2_1.png}
    \caption{Actuator force vs.\ relative velocity, for the three
    control strategies at $V=10$, $30$, $50$ m/s.}
    \label{fig:ex2_1}
\end{figure}

In Figure~\ref{fig:ex2_1} the passive damper produces a thin line
with slope $c=1500$\,Ns/m, since $F = c\,v_{rel}$. The skyhook and
groundhook clouds are noticeably broader: the same relative velocity
can correspond to widely different forces depending on the absolute
sprung and unsprung velocities. As the speed increases the road
excitation becomes more energetic, and both the relative velocity
range and the actuator force span grow accordingly. This confirms
that the actuator must deliver larger forces at higher speeds,
information that is critical for component sizing.

%======================================================================
\section{Exercise 2.2 -- ``Real-World'' Skyhook on a Semi-Active Suspension}
%======================================================================

\subsection{Theory}
A semi-active damper can only dissipate energy: the force must be
opposite to the relative velocity. The \emph{real-world skyhook}
control law selects between an active-like force and a purely passive
one according to a sign condition:
\begin{equation}
    F =
    \begin{cases}
        -(c_s + c)\dot z_s + c\,\dot z_u, &
            \text{if } \dot z_s(\dot z_s - \dot z_u) > 0 ,\\[2pt]
        -c(\dot z_s - \dot z_u), & \text{otherwise.}
    \end{cases}
    \label{eq:rw_sky}
\end{equation}
The first branch is active when the chassis is moving in the same
direction as the damper would naturally push it; the second branch
keeps the damper passive otherwise. The constraint
$c\,v_{rel} \le F \le c_{max}\,v_{rel}$ is then imposed by a
saturation, with $c_{max} = 4000$\,Ns/m.

\subsection{Implementation}

\begin{figure}[H]
    \centering
    \includegraphics[width=0.95\textwidth]{images/sim2_2.png}
    \caption{Simulink model \texttt{exercise2\_2SIM.slx}: the
    real-world skyhook law is implemented inside the
    \emph{MATLAB Function} block and applied as a second input to the
    state-space, after a 100 Hz low-pass filter that prevents
    chattering.}
    \label{fig:sim2_2}
\end{figure}

In the model of Figure~\ref{fig:sim2_2} the state-space block uses
the passive matrices (no damping inside $A$), while the actuator
force $F$ enters $B$ as a separate input. A MATLAB Function block
implements~\eqref{eq:rw_sky} using the sprung and unsprung velocities
fed back from the state-space. A first-order low-pass filter with cut
frequency $100$\,Hz smooths the discontinuity introduced by the
switching law. The MATLAB driver sweeps the tuning parameter $c_s$:
\begin{lstlisting}
cs_tune = [0, 500, 1000, 2000, 3000, 4000];
for i = 1:length(cs_tune)
    cs = cs_tune(i);
    out = sim('exercise2_2SIM.slx');
    rms_acc(i) = sqrt(mean(out.acc.^2));
end
\end{lstlisting}

\subsection{Results}

\begin{figure}[H]
    \centering
    \includegraphics[width=0.95\textwidth]{images/ex2_2.png}
    \caption{Sprung-mass acceleration time history for $c_s = 0$
    (passive) and five semi-active configurations.}
    \label{fig:ex2_2}
\end{figure}

\begin{figure}[H]
    \centering
    \includegraphics[width=0.7\textwidth]{images/ex2_2-results.png}
    \caption{RMS sprung-mass acceleration and percentage reduction
    versus passive case.}
    \label{fig:ex2_2_results}
\end{figure}

Figures~\ref{fig:ex2_2} and~\ref{fig:ex2_2_results} show that even a
modest semi-active correction reduces the RMS acceleration: from
$0.72$ m/s\textsuperscript{2} (passive, $c_s = 0$) to a minimum of
$0.66$ m/s\textsuperscript{2} at $c_s = 2000$\,Ns/m, that is a
\textbf{8.2\%} comfort improvement. Increasing $c_s$ further is
counter-productive: the curves at $c_s = 3000$ and $4000$\,Ns/m show
slightly higher RMS values, because the saturation activates more
often and the switching law spends more time in the passive branch.
The optimum is therefore $c_s \approx 2000$\,Ns/m, which sits roughly
half-way between the passive coefficient and the prescribed
$c_{max}$.

%======================================================================
\section{Exercise 3 (Part 1) -- Nonlinear Semi-Active Suspension}
%======================================================================

\subsection{Theory}
A magneto-rheological (MR) damper provides a nonlinear force
$u(I, v_{wc})$ as a function of the control current $I$ and the
piston velocity $v_{wc} = \dot z_u - \dot z_s$. The nonlinear
quarter-car equations are
\begin{align}
    m_u \ddot z_u &= -u(I, v_{wc}) - K(z_u - z_s) - P(z_u - r) -
                     c_t(\dot z_u - \dot r),\\
    m_s \ddot z_s &= u(I, v_{wc}) + K(z_u - z_s).
\end{align}
The outputs are the chassis vertical acceleration
$y_1 = \tfrac{c}{m_s}(\dot z_u - \dot z_s) + \tfrac{K}{m_s}(z_u -
z_s)$ and the dynamic tyre-road force
$y_2 = P(z_u - r)$. A proportional feedback law manipulates the
current,
\begin{equation}
    I = \mathrm{sat}\bigl(I_c - K_p\,y_1^d\bigr) ,
    \quad I \in [0.29,\,1.6]\,\text{A},
    \label{eq:ex3_ctrl}
\end{equation}
where $y_1^d$ is $y_1$ delayed by a Unit Delay block with
$\tau=0.005$\,s, $I_c$ is the upper-bound current and $K_p$ is the
tuning gain. The closed-loop adapts the damping in real time: when the
chassis accelerates upward, the current is reduced, the damper softens
and the disturbance is filtered; when the chassis accelerates
downward, the current is increased and the damper hardens.

\subsection{Implementation}

\begin{figure}[H]
    \centering
    \includegraphics[width=\textwidth]{images/sim3.png}
    \caption{Simulink model \texttt{exercise\_3sim.slx}. From left to
    right: random/bump selector, 2-DOF nonlinear quarter-car (four
    integrators), 2-D lookup table for the MR damper, current control
    loop with Unit Delay, gain $K_p$ and Saturation.}
    \label{fig:sim3}
\end{figure}

The model in Figure~\ref{fig:sim3} contains four integrators for the
states $\dot z_s, z_s, \dot z_u, z_u$, the algebraic relations for
$y_1$ and $y_2$, the 2-D lookup table that returns the actuator
force, and the discrete control branch on the right. The selector at
the input switches between a random profile (band-limited noise
shaped as in Exercise 2.1) and a bump shape, controlled by the
workspace variable \texttt{inp}.

\paragraph{$K_p$ tuning at $I_c = 1.6$\,A.}
The proportional gain is swept in $K_p \in [0.2, 0.6]$ and four
performance indicators are recorded for each road profile: RMS and
maximum chassis acceleration, RMS and maximum tyre-ground force.

\begin{lstlisting}
K_p_tuning = linspace(0.2,0.6,20);
for k = 1:2          % 1 = random, 2 = bump
    inp = k;
    for i = 1:length(K_p_tuning)
        K_p = K_p_tuning(i);
        sim("exercise_3sim.slx");
        acc_RMS(i)   = sqrt(mean(Acc_sprung.^2));
        force_RMS(i) = sqrt(mean(Tire_road_force.^2));
        acc_Max(i)   = max(Acc_sprung);
        force_Max(i) = max(Tire_road_force);
    end
end
\end{lstlisting}

\subsection{Results -- $K_p$ tuning}

\begin{figure}[H]
    \centering
    \includegraphics[width=\textwidth]{images/ex3_kp_rand.png}
    \caption{$K_p$ optimisation for a random road profile,
    $I_c = 1.6$\,A.}
    \label{fig:ex3_kp_rand}
\end{figure}

\begin{figure}[H]
    \centering
    \includegraphics[width=\textwidth]{images/ex3_kp_bump.png}
    \caption{$K_p$ optimisation for a bump road profile,
    $I_c = 1.6$\,A.}
    \label{fig:ex3_kp_bump}
\end{figure}

For the \textbf{random profile}
(Figure~\ref{fig:ex3_kp_rand}) the four metrics point to different
optima:
\begin{itemize}
    \item RMS acceleration minimised at $K_p = 0.33$
          ($1.04$ m/s\textsuperscript{2});
    \item RMS tyre force monotonically decreasing, lowest at
          $K_p = 0.58$ ($\approx 440$ N);
    \item maximum acceleration minimum at $K_p = 0.26$
          ($2.69$ m/s\textsuperscript{2}), then jumps up sharply for
          $K_p \ge 0.3$;
    \item maximum tyre force monotonically decreasing.
\end{itemize}
The classic \emph{comfort-vs-handling} trade-off emerges: comfort
metrics favour low $K_p$ (soft damper, large body excursion absorbed
by the suspension), while handling metrics prefer high $K_p$ (firm
damper, smaller wheel hop).

For the \textbf{bump profile}
(Figure~\ref{fig:ex3_kp_bump}) the four optima are much closer:
$K_p \approx 0.24$--$0.31$ minimises both RMS metrics; the maximum
tyre force has its lowest value at $K_p \approx 0.41$. The transient
nature of a bump excitation aligns comfort and handling indicators on
a common interval. A reasonable single value covering both scenarios
is $K_p \approx 0.30$: near-optimal on the bump and close to the
RMS-acceleration optimum on the random profile, accepting a sub-optimal
but still acceptable RMS force.

\subsection{Results -- $I_c$ tuning at fixed $K_p = 0.3$}

\begin{figure}[H]
    \centering
    \includegraphics[width=\textwidth]{images/ex3_Ic_rand.png}
    \caption{$I_c$ optimisation for a random road profile,
    $K_p = 0.3$.}
    \label{fig:ex3_Ic_rand}
\end{figure}

\begin{figure}[H]
    \centering
    \includegraphics[width=\textwidth]{images/ex3_ic_bump.png}
    \caption{$I_c$ optimisation for a bump road profile,
    $K_p = 0.3$.}
    \label{fig:ex3_ic_bump}
\end{figure}

For the \textbf{random profile}
(Figure~\ref{fig:ex3_Ic_rand}) the RMS acceleration grows almost
linearly with $I_c$, from $0.53$ to $1.05$ m/s\textsuperscript{2}: a
softer baseline damper is better for comfort. The RMS tyre force has
a clear minimum at $I_c \approx 0.5$\,A; below this point the damper
is too compliant and the wheel oscillates excessively, above it the
sprung-mass motion injects more energy into the tyre. The maximum
indicators show a similar non-monotonic behaviour with a minimum
around $0.5$\,A.

For the \textbf{bump profile}
(Figure~\ref{fig:ex3_ic_bump}) the picture changes: the RMS
acceleration is minimised around $I_c \approx 1.0$\,A, and so is the
RMS tyre force. A bump is a deterministic, low-energy impulse, and
extra damping helps to absorb it quickly without amplifying tyre
hopping.

\paragraph{Take-away.}
The optimal current depends strongly on the type of excitation. A
practical implementation therefore requires either a road-profile
classifier (e.g., a low-pass-filtered estimator of road harshness)
or, alternatively, a robust compromise such as $I_c \approx 0.5$\,A
on random pavement and $I_c \approx 1.0$\,A in transient manoeuvres.

%======================================================================
\section{Exercise 3 (Part 2) -- Bode of the Linearised System}
%======================================================================

\subsection{Theory}
The nonlinear system can be linearised in Simulink through the
\emph{Linear Analysis Tool}: input/output perturbation points are
added on the road derivative and on the two outputs $y_1, y_2$;
Simulink computes a Jacobian linearisation around the chosen
operating point. The resulting state-space approximates the closed
loop with the lookup table replaced by its local slope
$\partial u / \partial v_{wc}$, that is, an equivalent linear damper
of coefficient $c_{eq}(I)$. As $I$ varies between $0.29$\,A and
$1.6$\,A the equivalent damping spans more than one order of
magnitude.

\subsection{Implementation}
The linearisation session was saved as \texttt{linearization.mat},
which contains two \texttt{ss} objects named \texttt{Low\_current}
($I_c = 0.29$\,A) and \texttt{High\_current} ($I_c = 1.6$\,A). They
are reloaded and used to draw a Bode plot directly from MATLAB:
\begin{lstlisting}
S = load('linearization.mat');
for k = 1:numel(S.LocalVariables)
    switch S.LocalVariables(k).Name
        case 'Low_current',  sys_low  = S.LocalVariables(k).Value;
        case 'High_current', sys_high = S.LocalVariables(k).Value;
    end
end
figure
bode(sys_low, 'b', sys_high, 'r');
legend('I_c = 0.29 A','I_c = 1.6 A','Location','best')
grid on
\end{lstlisting}

\subsection{Results}

\begin{figure}[H]
    \centering
    \includegraphics[width=0.85\textwidth]{images/ex3_bode.png}
    \caption{Bode diagram of the linearised semi-active system. Top:
    transfer function to the tyre-road force; bottom: transfer
    function to the chassis vertical acceleration.}
    \label{fig:ex3_bode}
\end{figure}

The Bode plots in Figure~\ref{fig:ex3_bode} reproduce the two-mode
structure already observed in Exercise 1: the chassis resonance
appears around $7$\,rad/s and the wheel resonance around
$55$--$60$\,rad/s. The high-current condition shows a more damped
chassis peak (better comfort at low frequencies) but a slightly more
pronounced wheel resonance, while the low-current condition produces
the opposite. This is exactly the ``skyhook envelope'' that the
proportional controller \eqref{eq:ex3_ctrl} aims to exploit: by
moving along the current axis in real time, the closed loop can
choose, frequency by frequency, the most favourable response between
the two extremes. Compared to the linear analysis of Exercise 1, the
peaks here are slightly broader because the linearisation absorbs
the nonlinear lookup at a single operating point, but the overall
agreement is very satisfactory.

%======================================================================
\section{Conclusions}
%======================================================================
The project provided a complete view of vertical vehicle dynamics
seen from the angle of suspension control, moving from idealised
linear analysis to a realistic semi-active implementation with a
magneto-rheological damper.

\begin{itemize}
    \item Frequency-domain analysis (Exercise 1) confirmed the
          complementary nature of skyhook and groundhook: the former
          excels around the chassis mode, the latter around the wheel
          mode. Eigenvalue analysis showed how the two strategies
          relocate the corresponding pole pairs.
    \item Time-domain simulation (Exercise 2.1) translated these
          observations into actuator-force/relative-velocity
          diagrams, useful for component sizing across the operating
          speed range.
    \item The real-world skyhook on a semi-active suspension
          (Exercise 2.2) reduced the RMS chassis acceleration by
          $\sim 8\%$ with a moderate tuning ($c_s \approx 2000$\,Ns/m),
          while respecting the dissipative constraint.
    \item The nonlinear MR-damper model (Exercise 3) showed that the
          control gain $K_p$ and the baseline current $I_c$ must be
          chosen jointly with the type of road excitation. A
          compromise tuning $K_p \approx 0.3$, $I_c \approx
          0.5$--$1.0$\,A behaved consistently on both bump and random
          profiles.
    \item The Bode analysis of the linearised closed loop
          (Exercise 3, Part 2) closed the picture by confirming that
          the semi-active controller spans an envelope between the
          two extreme passive equivalents, formally justifying the
          observed performance.
\end{itemize}

Overall, the model has proved consistent across all three exercises:
linear predictions, time-domain simulations and linearised Bode of
the nonlinear system all match within the expected accuracy.
